\begin{solapartat}
\punts{0,15} El so que arriba a cada persona és la interferència dels dos sons emesos. La intensitat serà màxima en el punt A quan tinguem interferència constructiva, és a dir quan la longitud d'ona i la diferència de camins tinguin una relació entera: $|d_2 - d_1| = 3,68 = n\lambda$ amb n enter.

\punts{0,10} les freqüències corresponents seran: $f = \frac{v}{\lambda} = n\,\frac{340}{3,68}\ Hz$.

\punts{0,25} En el rang de 200 a 400 Hz aquestes freqüències són:

$277,17\,Hz$;\ \ $369,56\ Hz$;

\punts{0,15} Perquè la persona B no pugui escoltar el senyal (l'escolti amb intensitat mínima) hem de tenir interferència destructiva en el punt B. En aquest cas la longitud d'ona i la diferència de camins han de tenir una relació semientera: $|d'_2 - d'_1| = 4,14 = \left(n + \frac{1}{2}\right)\lambda$ amb n enter

\punts{0,10} les freqüències corresponents seran: $f = \frac{v}{\lambda} = \left(n + \frac{1}{2}\right)\frac{340}{4,14}\ Hz$.

\punts{0,5} En el rang de 200 a 400 Hz aquestes freqüències són:

$205,31\ Hz$;\ \ $287,44\ Hz$;\ \ $369,56\ Hz$;

Observem que $f = 369,56\ Hz$ compleix les dues condicions.
\end{solapartat}

\begin{solapartat}
\punts{0,5} A partir de la potència emesa i les distàncies, calculem la intensitat sonora generada per cada un dels altaveus que arriba al punt A:
\begin{align*}
I_1 &= \frac{P}{4\pi d_1^2} = \frac{2,36\cdot 10^{-3}}{4\pi\cdot 5^2} = 7,51\cdot 10^{-6}\ W/m^2\\
I_2 &= \frac{P}{4\pi d_2^2} = \frac{2,36\cdot 10^{-3}}{4\pi\cdot 8,68^2} = 2,49\cdot 10^{-6}\ W/m^2
\end{align*}

\punts{0,25} La intensitat total és la suma de les intensitats que ens arriben de cada altaveu: $I = I_1 + I_2 = 7,51\cdot 10^{-6} + 2,49\cdot 10^{-6} = 10^{-5}\ W/m^2$

\punts{0,5} El nivell d'intensitat sonora percebut per la persona que es troba al punt A és:
\[ \beta = 10\log\frac{I}{I_0} = 10\log\frac{10^{-5}}{10^{-12}} = 70\ dB\,. \]
\end{solapartat}
